% Generated. Do not edit.
% Projection scriptural-date-location/section, version 1: the Scriptural Date and
% Location page, generated whole from research/chronology.toml, the Scripture
% chronology corpus and the source library by tools/tpt document-types refresh.
% document: liturgy/roman-rite/1962/propers/temporal/66-resumed-fifth-after-epiphany
\runninghead{Scriptural Date and Location}
\section{Scriptural Date and Location}
\label{sec:date-location}

{\small Each passage the Mass reads or sings from Scripture, in canonical order. The Context column names what is dated --- the event the passage narrates, a setting it remembers, the traditional attribution of its author, or the composition of its book --- and where it belongs; the Date column gives the answer of the Scripture chronology corpus for that subject, with every alternative it records. Texts of the Mass that are not drawn from Scripture have no row. The evidence for each date follows the table.\par}
\medskip
% Generated. Do not edit.
% Deterministic projection of research/chronology.toml.
% schema: 5
% document: liturgy/roman-rite/1962/propers/temporal/66-resumed-fifth-after-epiphany
% generated-by: tools/tpt proper-chronology annotations
\newcommand{\chronologyannotationclaim}[6]{#6}
\newcommand{\chronologyannotationreach}[3]{}
\newcommand{\chronologyannotationgroup}[3]{#3}
\newcommand{\chronologyannotation}[1]{%
  \ifcsname triptychchronologyannotation@#1\endcsname
    \csname triptychchronologyannotation@#1\endcsname
  \else
    \PackageError{triptych}{No chronology annotation for #1}{}%
  \fi}
% introit: Introit
\expandafter\def\csname triptychchronologyannotation@introit\endcsname{%
  \chronologyannotationgroup{historical-setting}{research-pending}{\textbf{Historical setting} -- Occasion date unresolved.}%
  \space \chronologyannotationgroup{composition}{nonuniform}{\textbf{Composition} -- No single date applies across every cited locus.}%
}
% epistle: Epistle
\expandafter\def\csname triptychchronologyannotation@epistle\endcsname{%
  \chronologyannotationgroup{composition}{disputed}{\textbf{Composition} -- disputed: \chronologyannotationreach{Col.3.12}{Col}{true}\chronologyannotationreach{Col.3.13}{Col}{true}\chronologyannotationreach{Col.3.14}{Col}{true}\chronologyannotationreach{Col.3.15}{Col}{true}\chronologyannotationreach{Col.3.16}{Col}{true}\chronologyannotationreach{Col.3.17}{Col}{true}\chronologyannotationclaim{composition.epistle-to-the-colossians}{composition}{catholic-traditional-v1}{disputed}{(Philemon; Colossians; Ephesians; Philippians), 61}{A.D. 61}; \chronologyannotationreach{Col.3.12}{Col}{true}\chronologyannotationreach{Col.3.13}{Col}{true}\chronologyannotationreach{Col.3.14}{Col}{true}\chronologyannotationreach{Col.3.15}{Col}{true}\chronologyannotationreach{Col.3.16}{Col}{true}\chronologyannotationreach{Col.3.17}{Col}{true}\chronologyannotationclaim{composition.epistle-to-the-colossians}{composition}{catholic-traditional-v1}{disputed}{some time between A.D. 61 and 63}{A.D. 61--63}.}%
}
% gradual: Gradual
\expandafter\def\csname triptychchronologyannotation@gradual\endcsname{%
  \chronologyannotationgroup{historical-setting}{research-pending}{\textbf{Historical setting} -- Occasion date unresolved.}%
  \space \chronologyannotationgroup{composition}{preferred}{\textbf{Composition}: \chronologyannotationreach{Ps.43.8}{Ps.41 Ps.42 Ps.43 Ps.44 Ps.45 Ps.46 Ps.47 Ps.48 Ps.83 Ps.84 Ps.86 Ps.87}{true}\chronologyannotationreach{Ps.43.9}{Ps.41 Ps.42 Ps.43 Ps.44 Ps.45 Ps.46 Ps.47 Ps.48 Ps.83 Ps.84 Ps.86 Ps.87}{true}\chronologyannotationclaim{composition.psalms-of-the-sons-of-korah}{composition}{catholic-traditional-v1}{preferred}{between the days of Isaias and the return from exile}{Between the days of Isaias and the return from exile}.}%
}
% alleluia: Alleluia
\expandafter\def\csname triptychchronologyannotation@alleluia\endcsname{%
  \chronologyannotationgroup{historical-setting}{research-pending}{\textbf{Historical setting} -- Occasion date unresolved.}%
  \space \chronologyannotationgroup{composition}{preferred}{\textbf{Composition}: \chronologyannotationreach{Ps.129.1}{Ps}{true}\chronologyannotationreach{Ps.129.2}{Ps}{true}\chronologyannotationclaim{critical.psalms.latest-composition-boundary}{composition}{catholic-critical-v1}{preferred}{before the Maccabean period, around 165 B.C.; no individual psalm can be dated securely}{Before c. 165 B.C.}}%
}
% gospel: Gospel
\expandafter\def\csname triptychchronologyannotation@gospel\endcsname{%
  \chronologyannotationgroup{narrated-event}{research-pending}{\textbf{Event} -- No narrated-event date in the chronology corpus.}%
  \space \chronologyannotationgroup{composition}{disputed}{\textbf{Composition} -- disputed: \chronologyannotationreach{Matt.13.24}{Matt}{true}\chronologyannotationreach{Matt.13.25}{Matt}{true}\chronologyannotationreach{Matt.13.26}{Matt}{true}\chronologyannotationreach{Matt.13.27}{Matt}{true}\chronologyannotationreach{Matt.13.28}{Matt}{true}\chronologyannotationreach{Matt.13.29}{Matt}{true}\chronologyannotationreach{Matt.13.30}{Matt}{true}\chronologyannotationclaim{composition.gospel-of-matthew}{composition}{catholic-traditional-v1}{disputed}{about A.D. 38-45}{c. A.D. 38--45}; \chronologyannotationreach{Matt.13.24}{Matt}{true}\chronologyannotationreach{Matt.13.25}{Matt}{true}\chronologyannotationreach{Matt.13.26}{Matt}{true}\chronologyannotationreach{Matt.13.27}{Matt}{true}\chronologyannotationreach{Matt.13.28}{Matt}{true}\chronologyannotationreach{Matt.13.29}{Matt}{true}\chronologyannotationreach{Matt.13.30}{Matt}{true}\chronologyannotationclaim{composition.gospel-of-matthew}{composition}{catholic-traditional-v1}{disputed}{about the year 40-42}{c. A.D. 40--42}; \chronologyannotationreach{Matt.13.24}{Matt}{true}\chronologyannotationreach{Matt.13.25}{Matt}{true}\chronologyannotationreach{Matt.13.26}{Matt}{true}\chronologyannotationreach{Matt.13.27}{Matt}{true}\chronologyannotationreach{Matt.13.28}{Matt}{true}\chronologyannotationreach{Matt.13.29}{Matt}{true}\chronologyannotationreach{Matt.13.30}{Matt}{true}\chronologyannotationclaim{composition.gospel-of-matthew}{composition}{catholic-traditional-v1}{disputed}{the years 40-45}{A.D. 40--45}; \chronologyannotationreach{Matt.13.24}{Matt}{true}\chronologyannotationreach{Matt.13.25}{Matt}{true}\chronologyannotationreach{Matt.13.26}{Matt}{true}\chronologyannotationreach{Matt.13.27}{Matt}{true}\chronologyannotationreach{Matt.13.28}{Matt}{true}\chronologyannotationreach{Matt.13.29}{Matt}{true}\chronologyannotationreach{Matt.13.30}{Matt}{true}\chronologyannotationclaim{composition.gospel-of-matthew}{composition}{catholic-traditional-v1}{disputed}{about the year 60-68}{c. A.D. 60--68}; \chronologyannotationreach{Matt.13.24}{Matt}{true}\chronologyannotationreach{Matt.13.25}{Matt}{true}\chronologyannotationreach{Matt.13.26}{Matt}{true}\chronologyannotationreach{Matt.13.27}{Matt}{true}\chronologyannotationreach{Matt.13.28}{Matt}{true}\chronologyannotationreach{Matt.13.29}{Matt}{true}\chronologyannotationreach{Matt.13.30}{Matt}{true}\chronologyannotationclaim{composition.gospel-of-matthew}{composition}{catholic-traditional-v1}{disputed}{about the years 64-67}{c. A.D. 64--67}; \chronologyannotationreach{Matt.13.24}{Matt}{true}\chronologyannotationreach{Matt.13.25}{Matt}{true}\chronologyannotationreach{Matt.13.26}{Matt}{true}\chronologyannotationreach{Matt.13.27}{Matt}{true}\chronologyannotationreach{Matt.13.28}{Matt}{true}\chronologyannotationreach{Matt.13.29}{Matt}{true}\chronologyannotationreach{Matt.13.30}{Matt}{true}\chronologyannotationclaim{composition.gospel-of-matthew}{composition}{catholic-traditional-v1}{disputed}{about the year 50}{c. A.D. 50}.}%
}
% offertory: Offertory
\expandafter\def\csname triptychchronologyannotation@offertory\endcsname{%
  \chronologyannotationgroup{historical-setting}{research-pending}{\textbf{Historical setting} -- Occasion date unresolved.}%
  \space \chronologyannotationgroup{composition}{preferred}{\textbf{Composition}: \chronologyannotationreach{Ps.129.1}{Ps}{true}\chronologyannotationreach{Ps.129.2}{Ps}{true}\chronologyannotationclaim{critical.psalms.latest-composition-boundary}{composition}{catholic-critical-v1}{preferred}{before the Maccabean period, around 165 B.C.; no individual psalm can be dated securely}{Before c. 165 B.C.}}%
}
% communion: Communion
\expandafter\def\csname triptychchronologyannotation@communion\endcsname{%
  \chronologyannotationgroup{composition}{disputed}{\textbf{Composition} -- disputed: \chronologyannotationreach{Mark.11.24}{Mark}{true}\chronologyannotationclaim{composition.gospel-of-mark}{composition}{catholic-traditional-v1}{disputed}{the tenth or twelfth year after the Ascension (A.D. 38-40)}{A.D. 38--40}; \chronologyannotationreach{Mark.11.24}{Mark}{true}\chronologyannotationclaim{composition.gospel-of-mark}{composition}{catholic-traditional-v1}{disputed}{A.D. 40}{A.D. 40}; \chronologyannotationreach{Mark.11.24}{Mark}{true}\chronologyannotationclaim{composition.gospel-of-mark}{composition}{catholic-traditional-v1}{disputed}{the third year of Claudius (A.D. 43)}{A.D. 43}; \chronologyannotationreach{Mark.11.24}{Mark}{true}\chronologyannotationclaim{composition.gospel-of-mark}{composition}{catholic-traditional-v1}{disputed}{somewhere between A.D. 50 and 67}{A.D. 50--67}; \chronologyannotationreach{Mark.11.24}{Mark}{true}\chronologyannotationclaim{composition.gospel-of-mark}{composition}{catholic-traditional-v1}{disputed}{between 60 and 70}{A.D. 60--70}; \chronologyannotationreach{Mark.11.24}{Mark}{true}\chronologyannotationclaim{composition.gospel-of-mark}{composition}{catholic-traditional-v1}{disputed}{after Matthew, and not as late as the destruction of Jerusalem}{After Matthew, and not as late as the destruction of Jerusalem}.}%
}
% presentation: orientation-v3
\newcommand{\chronologyorientationrow}[2]{%
  \ifcsname triptychchronologyorientation@#1\endcsname
    \csname triptychchronologyorientation@#1\endcsname{#2}%
  \else
    \PackageError{triptych}{No chronology orientation row for #1}{}%
  \fi}
\expandafter\def\csname triptychchronologyorientation@row-a16db10bcf3bfa76\endcsname#1{%
Gradual: Ps. 43:8--9 (Vg.; Heb. 44) & Historical setting & Occasion date unresolved\\*[.1em]%
}
\expandafter\def\csname triptychchronologyorientation@row-fd3920d9cd8b2ce1\endcsname#1{%
 & Composition: The Psalms of the Sons of Korah & \chronologyannotationreach{Ps.43.8}{Ps.41 Ps.42 Ps.43 Ps.44 Ps.45 Ps.46 Ps.47 Ps.48 Ps.83 Ps.84 Ps.86 Ps.87}{true}\chronologyannotationreach{Ps.43.9}{Ps.41 Ps.42 Ps.43 Ps.44 Ps.45 Ps.46 Ps.47 Ps.48 Ps.83 Ps.84 Ps.86 Ps.87}{true}\chronologyannotationclaim{composition.psalms-of-the-sons-of-korah}{composition}{catholic-traditional-v1}{preferred}{between the days of Isaias and the return from exile}{Between the days of Isaias and the return from exile}\\[.3em]%
}
\expandafter\def\csname triptychchronologyorientation@row-daf59e02c6c15663\endcsname#1{%
Introit: Ps. 84:2 (Vg.; Heb. 85) & Historical setting & Occasion date unresolved\\*[.1em]%
}
\expandafter\def\csname triptychchronologyorientation@row-d85392b5260231db\endcsname#1{%
 & Composition: The Psalms of the Sons of Korah & \chronologyannotationreach{Ps.84.2}{Ps.41 Ps.42 Ps.43 Ps.44 Ps.45 Ps.46 Ps.47 Ps.48 Ps.83 Ps.84 Ps.86 Ps.87}{true}\chronologyannotationclaim{composition.psalms-of-the-sons-of-korah}{composition}{catholic-traditional-v1}{preferred}{between the days of Isaias and the return from exile}{Between the days of Isaias and the return from exile}\\[.3em]%
}
\expandafter\def\csname triptychchronologyorientation@row-e61fc8dcfe43e661\endcsname#1{%
Alleluia: Ps. 129:1--2 (Vg.; Heb. 130) & Historical setting & Occasion date unresolved\\*[.1em]%
}
\expandafter\def\csname triptychchronologyorientation@row-777243e0b53190a9\endcsname#1{%
 & Composition: Psalms & \chronologyannotationreach{Ps.129.1}{Ps}{true}\chronologyannotationreach{Ps.129.2}{Ps}{true}\chronologyannotationclaim{critical.psalms.latest-composition-boundary}{composition}{catholic-critical-v1}{preferred}{before the Maccabean period, around 165 B.C.; no individual psalm can be dated securely}{Before c.~165 B.C.}\\[.3em]%
}
\expandafter\def\csname triptychchronologyorientation@row-33a9b5f385b6bb7a\endcsname#1{%
Offertory: Ps. 129:1--2 (Vg.; Heb. 130) & Historical setting & Occasion date unresolved\\*[.1em]%
}
\expandafter\def\csname triptychchronologyorientation@row-42f28ec9d138025a\endcsname#1{%
 & Composition: Psalms & \chronologyannotationreach{Ps.129.1}{Ps}{true}\chronologyannotationreach{Ps.129.2}{Ps}{true}\chronologyannotationclaim{critical.psalms.latest-composition-boundary}{composition}{catholic-critical-v1}{preferred}{before the Maccabean period, around 165 B.C.; no individual psalm can be dated securely}{Before c.~165 B.C.}\\[.3em]%
}
\expandafter\def\csname triptychchronologyorientation@row-ea8e2406b53ffb05\endcsname#1{%
Introit: Jer. 29:11--12,14 & Composition: Jeremias's letter to the exiles after Jechonias's deportation; Jerusalem to Babylon & \chronologyannotationreach{Jer.29.11}{Jer.29.1-23}{false}\chronologyannotationreach{Jer.29.12}{Jer.29.1-23}{false}\chronologyannotationreach{Jer.29.14}{Jer.29.1-23}{false}\chronologyannotationclaim{composition.jeremias-letter-to-the-exiles}{composition}{catholic-traditional-v1}{preferred}{After that Jechonias the king ... were departed out of Jerusalem}{After Jechonias and the deportees had departed from Jerusalem}\\*[.1em]%
}
\expandafter\def\csname triptychchronologyorientation@row-956fb69c94da31fb\endcsname#1{%
 & Final formation: The received Book of Jeremias after enlargement; Judah ministry; final writing site unspecified & \chronologyannotationreach{Jer.29.11}{Jer}{true}\chronologyannotationreach{Jer.29.12}{Jer}{true}\chronologyannotationreach{Jer.29.14}{Jer}{true}\chronologyannotationclaim{composition.book-of-jeremias}{final-formation}{catholic-traditional-v1}{preferred}{the first transcript was enlarged by various insertions and additions}{After the first authentic recension}\\[.3em]%
}
\expandafter\def\csname triptychchronologyorientation@row-589914bfd2c486c9\endcsname#1{%
Gospel: Matt. 13:24--30 & Narrated event: Jesus tells the parable of the wheat and cockle beside the sea & Date unresolved\\*[.1em]%
}
\expandafter\def\csname triptychchronologyorientation@row-a387763c6e967d19\endcsname#1{%
 & Composition: Matthew; Semitic original associated with Palestine & Disputed: \chronologyannotationreach{Matt.13.24}{Matt}{true}\chronologyannotationreach{Matt.13.25}{Matt}{true}\chronologyannotationreach{Matt.13.26}{Matt}{true}\chronologyannotationreach{Matt.13.27}{Matt}{true}\chronologyannotationreach{Matt.13.28}{Matt}{true}\chronologyannotationreach{Matt.13.29}{Matt}{true}\chronologyannotationreach{Matt.13.30}{Matt}{true}\chronologyannotationclaim{composition.gospel-of-matthew}{composition}{catholic-traditional-v1}{disputed}{about A.D. 38-45}{c.~A.D.~38--45}; \chronologyannotationreach{Matt.13.24}{Matt}{true}\chronologyannotationreach{Matt.13.25}{Matt}{true}\chronologyannotationreach{Matt.13.26}{Matt}{true}\chronologyannotationreach{Matt.13.27}{Matt}{true}\chronologyannotationreach{Matt.13.28}{Matt}{true}\chronologyannotationreach{Matt.13.29}{Matt}{true}\chronologyannotationreach{Matt.13.30}{Matt}{true}\chronologyannotationclaim{composition.gospel-of-matthew}{composition}{catholic-traditional-v1}{disputed}{about the year 40-42}{c.~A.D.~40--42}; \chronologyannotationreach{Matt.13.24}{Matt}{true}\chronologyannotationreach{Matt.13.25}{Matt}{true}\chronologyannotationreach{Matt.13.26}{Matt}{true}\chronologyannotationreach{Matt.13.27}{Matt}{true}\chronologyannotationreach{Matt.13.28}{Matt}{true}\chronologyannotationreach{Matt.13.29}{Matt}{true}\chronologyannotationreach{Matt.13.30}{Matt}{true}\chronologyannotationclaim{composition.gospel-of-matthew}{composition}{catholic-traditional-v1}{disputed}{the years 40-45}{A.D.~40--45}; \chronologyannotationreach{Matt.13.24}{Matt}{true}\chronologyannotationreach{Matt.13.25}{Matt}{true}\chronologyannotationreach{Matt.13.26}{Matt}{true}\chronologyannotationreach{Matt.13.27}{Matt}{true}\chronologyannotationreach{Matt.13.28}{Matt}{true}\chronologyannotationreach{Matt.13.29}{Matt}{true}\chronologyannotationreach{Matt.13.30}{Matt}{true}\chronologyannotationclaim{composition.gospel-of-matthew}{composition}{catholic-traditional-v1}{disputed}{about the year 60-68}{c.~A.D.~60--68}; \chronologyannotationreach{Matt.13.24}{Matt}{true}\chronologyannotationreach{Matt.13.25}{Matt}{true}\chronologyannotationreach{Matt.13.26}{Matt}{true}\chronologyannotationreach{Matt.13.27}{Matt}{true}\chronologyannotationreach{Matt.13.28}{Matt}{true}\chronologyannotationreach{Matt.13.29}{Matt}{true}\chronologyannotationreach{Matt.13.30}{Matt}{true}\chronologyannotationclaim{composition.gospel-of-matthew}{composition}{catholic-traditional-v1}{disputed}{about the years 64-67}{c.~A.D.~64--67}; \chronologyannotationreach{Matt.13.24}{Matt}{true}\chronologyannotationreach{Matt.13.25}{Matt}{true}\chronologyannotationreach{Matt.13.26}{Matt}{true}\chronologyannotationreach{Matt.13.27}{Matt}{true}\chronologyannotationreach{Matt.13.28}{Matt}{true}\chronologyannotationreach{Matt.13.29}{Matt}{true}\chronologyannotationreach{Matt.13.30}{Matt}{true}\chronologyannotationclaim{composition.gospel-of-matthew}{composition}{catholic-traditional-v1}{disputed}{about the year 50}{c.~A.D.~50}\\[.3em]%
}
\expandafter\def\csname triptychchronologyorientation@row-99decd09d9817db7\endcsname#1{%
Communion: Mark 11:24 & Composition: Mark; Rome traditionally; Egypt reported by Chrysostom & Disputed: \chronologyannotationreach{Mark.11.24}{Mark}{true}\chronologyannotationclaim{composition.gospel-of-mark}{composition}{catholic-traditional-v1}{disputed}{the tenth or twelfth year after the Ascension (A.D. 38-40)}{A.D.~38--40}; \chronologyannotationreach{Mark.11.24}{Mark}{true}\chronologyannotationclaim{composition.gospel-of-mark}{composition}{catholic-traditional-v1}{disputed}{A.D. 40}{A.D.~40}; \chronologyannotationreach{Mark.11.24}{Mark}{true}\chronologyannotationclaim{composition.gospel-of-mark}{composition}{catholic-traditional-v1}{disputed}{the third year of Claudius (A.D. 43)}{A.D.~43}; \chronologyannotationreach{Mark.11.24}{Mark}{true}\chronologyannotationclaim{composition.gospel-of-mark}{composition}{catholic-traditional-v1}{disputed}{somewhere between A.D. 50 and 67}{A.D.~50--67}; \chronologyannotationreach{Mark.11.24}{Mark}{true}\chronologyannotationclaim{composition.gospel-of-mark}{composition}{catholic-traditional-v1}{disputed}{between 60 and 70}{A.D.~60--70}; \chronologyannotationreach{Mark.11.24}{Mark}{true}\chronologyannotationclaim{composition.gospel-of-mark}{composition}{catholic-traditional-v1}{disputed}{after Matthew, and not as late as the destruction of Jerusalem}{After Matthew, and not as late as the destruction of Jerusalem}\\[.3em]%
}
\expandafter\def\csname triptychchronologyorientation@row-b5e644ac10786f4b\endcsname#1{%
Epistle: Col. 3:12--17 & Composition: Colossians; Rome traditionally; addressed to Colossae & Disputed: \chronologyannotationreach{Col.3.12}{Col}{true}\chronologyannotationreach{Col.3.13}{Col}{true}\chronologyannotationreach{Col.3.14}{Col}{true}\chronologyannotationreach{Col.3.15}{Col}{true}\chronologyannotationreach{Col.3.16}{Col}{true}\chronologyannotationreach{Col.3.17}{Col}{true}\chronologyannotationclaim{composition.epistle-to-the-colossians}{composition}{catholic-traditional-v1}{disputed}{(Philemon; Colossians; Ephesians; Philippians), 61}{A.D.~61}; \chronologyannotationreach{Col.3.12}{Col}{true}\chronologyannotationreach{Col.3.13}{Col}{true}\chronologyannotationreach{Col.3.14}{Col}{true}\chronologyannotationreach{Col.3.15}{Col}{true}\chronologyannotationreach{Col.3.16}{Col}{true}\chronologyannotationreach{Col.3.17}{Col}{true}\chronologyannotationclaim{composition.epistle-to-the-colossians}{composition}{catholic-traditional-v1}{disputed}{some time between A.D. 61 and 63}{A.D.~61--63}\\[.3em]%
}

% Generated. Do not edit.
% Complete orientation-v3 table; wording, dates and order belong to the corpus.
\begin{chronologysharedtable}
\chronologyorientationrow{row-a16db10bcf3bfa76}{}
\chronologyorientationrow{row-fd3920d9cd8b2ce1}{}
\midrule
\chronologyorientationrow{row-daf59e02c6c15663}{}
\chronologyorientationrow{row-d85392b5260231db}{}
\midrule
\chronologyorientationrow{row-e61fc8dcfe43e661}{}
\chronologyorientationrow{row-777243e0b53190a9}{}
\midrule
\chronologyorientationrow{row-33a9b5f385b6bb7a}{}
\chronologyorientationrow{row-42f28ec9d138025a}{}
\midrule
\chronologyorientationrow{row-ea8e2406b53ffb05}{}
\chronologyorientationrow{row-956fb69c94da31fb}{}
\midrule
\chronologyorientationrow{row-589914bfd2c486c9}{}
\chronologyorientationrow{row-a387763c6e967d19}{}
\midrule
\chronologyorientationrow{row-99decd09d9817db7}{}
\midrule
\chronologyorientationrow{row-b5e644ac10786f4b}{}
\end{chronologysharedtable}
{\footnotesize Psalm numbering: Vg. = Vulgate; Heb. = Hebrew.\par}


\subsection{Evidence for each date}
\label{sec:date-location-evidence}
{\small For each passage, as the Scriptural Date and Location table names it: what is dated, how the corpus weighs each answer, the source's own words, and the sources cited. A subject the corpus does not date is reported as such and given no date.\par}
{\small
\par\smallskip\noindent\textit{Gradual: Ps. 43:8--9 (Vg.; Heb. 44)}\par\nobreak
\begin{itemize}[leftmargin=1.2em,itemsep=1pt,topsep=1pt,parsep=0pt]
\item \textbf{Historical setting}. No date: the chronology corpus has not yet dated this (research pending).
\item \textbf{Composition} (preferred). The Psalms of the Sons of Korah: Between the days of Isaias and the return from exile (The Catholic Encyclopedia, Volume 12 (1911)).
\end{itemize}
\par\smallskip\noindent\textit{Introit: Ps. 84:2 (Vg.; Heb. 85)}\par\nobreak
\begin{itemize}[leftmargin=1.2em,itemsep=1pt,topsep=1pt,parsep=0pt]
\item \textbf{Historical setting}. As under Gradual: Ps. 43:8--9 (Vg.; Heb. 44).
\item \textbf{Composition}. As under Gradual: Ps. 43:8--9 (Vg.; Heb. 44).
\end{itemize}
\par\smallskip\noindent\textit{Alleluia: Ps. 129:1--2 (Vg.; Heb. 130)}\par\nobreak
\begin{itemize}[leftmargin=1.2em,itemsep=1pt,topsep=1pt,parsep=0pt]
\item \textbf{Historical setting}. As under Gradual: Ps. 43:8--9 (Vg.; Heb. 44).
\item \textbf{Composition} (preferred). Psalms: Before c.~165 B.C. (``before the Maccabean period, around 165 B.C.; no individual psalm can be dated securely''; Catholic critical chronology; New American Bible, Revised Edition --- USCCB web state (2026), Psalms introduction).
\end{itemize}
\par\smallskip\noindent\textit{Offertory: Ps. 129:1--2 (Vg.; Heb. 130)}\par\nobreak
\begin{itemize}[leftmargin=1.2em,itemsep=1pt,topsep=1pt,parsep=0pt]
\item \textbf{Historical setting}. As under Gradual: Ps. 43:8--9 (Vg.; Heb. 44).
\item \textbf{Composition}. As under Alleluia: Ps. 129:1--2 (Vg.; Heb. 130).
\end{itemize}
\par\smallskip\noindent\textit{Introit: Jer. 29:11--12,14}\par\nobreak
\begin{itemize}[leftmargin=1.2em,itemsep=1pt,topsep=1pt,parsep=0pt]
\item \textbf{Composition} (preferred). Jeremias's letter to the exiles after Jechonias's deportation: After Jechonias and the deportees had departed from Jerusalem (``After that Jechonias the king ... were departed out of Jerusalem''; The Holy Bible, Douay-Rheims (Challoner revision), Project Gutenberg eBook 1581: Jer. 29:1, 29:2, 29:3, 29:24). Place: Jerusalem to Babylon (The Holy Bible, Douay-Rheims (Challoner revision), Project Gutenberg eBook 1581: Jer. 29:1, 29:3).
\item \textbf{Final formation} (preferred). The received Book of Jeremias after enlargement: After the first authentic recension (``the first transcript was enlarged by various insertions and additions''; The Catholic Encyclopedia, Volume 8 (New York, 1910)). Place: Judah ministry; final writing site unspecified (The Catholic Encyclopedia, Volume 8 (New York, 1910)).
\end{itemize}
\par\smallskip\noindent\textit{Gospel: Matt. 13:24--30}\par\nobreak
\begin{itemize}[leftmargin=1.2em,itemsep=1pt,topsep=1pt,parsep=0pt]
\item \textbf{Narrated event}. Jesus tells the parable of the wheat and cockle beside the sea (The Holy Bible, Douay-Rheims (Challoner revision), Project Gutenberg eBook 1581: Matt. 13:1, 13:2, 13:24, 13:25, 13:30). No date: the chronology corpus has not yet dated this (research pending).
\item \textbf{Composition} (disputed). The Gospel of St Matthew: c.~A.D.~38--45 (``about A.D. 38-45''; The Catholic Encyclopedia, Volume 10 (1911); The Catholic Encyclopedia, Volume 14 (New York, 1912)); c.~A.D.~40--42 (``about the year 40-42''; The Catholic Encyclopedia, Volume 10 (1911)); A.D.~40--45 (``the years 40-45''; The Catholic Encyclopedia, Volume 10 (1911)); c.~A.D.~60--68 (``about the year 60-68''; The Catholic Encyclopedia, Volume 10 (1911)); c.~A.D.~64--67 (``about the years 64-67''; The Catholic Encyclopedia, Volume 10 (1911)); c.~A.D.~50 (``about the year 50''; The Catholic Encyclopedia, Volume 14 (New York, 1912)). Place: Semitic original associated with Palestine (The Catholic Encyclopedia, Volume 10 (1911)).
\end{itemize}
\par\smallskip\noindent\textit{Communion: Mark 11:24}\par\nobreak
\begin{itemize}[leftmargin=1.2em,itemsep=1pt,topsep=1pt,parsep=0pt]
\item \textbf{Composition} (disputed). The Gospel of St Mark: A.D.~38--40 (``the tenth or twelfth year after the Ascension (A.D. 38-40)''; The Catholic Encyclopedia, Volume 9 (1910)); A.D.~40 (The Catholic Encyclopedia, Volume 9 (1910)); A.D.~43 (``the third year of Claudius (A.D. 43)''; The Catholic Encyclopedia, Volume 9 (1910)); A.D.~50--67 (``somewhere between A.D. 50 and 67''; The Catholic Encyclopedia, Volume 9 (1910)); A.D.~60--70 (``between 60 and 70''; The Catholic Encyclopedia, Volume 14 (New York, 1912)); After Matthew, and not as late as the destruction of Jerusalem (De auctore, de tempore compositionis et de historica veritate Evangeliorum secundum Marcum et secundum Lucam (1912), responsa-i-ix). Place: Rome traditionally; Egypt reported by Chrysostom (The Catholic Encyclopedia, Volume 9 (1910)).
\end{itemize}
\par\smallskip\noindent\textit{Epistle: Col. 3:12--17}\par\nobreak
\begin{itemize}[leftmargin=1.2em,itemsep=1pt,topsep=1pt,parsep=0pt]
\item \textbf{Composition} (disputed). The Epistle to the Colossians: A.D.~61 (``(Philemon; Colossians; Ephesians; Philippians), 61''; The Catholic Encyclopedia, Volume 11 (New York, 1911)); A.D.~61--63 (``some time between A.D. 61 and 63''; The Catholic Encyclopedia, Volume 4 (New York, 1908)). Place: Rome traditionally; addressed to Colossae (The Catholic Encyclopedia, Volume 4 (New York, 1908); The Holy Bible, Douay-Rheims (Challoner revision), Project Gutenberg eBook 1581: Col. 1:2).
\end{itemize}
\par}
